\chapter{Stochastic Differential Equations and PDEs}\label{ch:sde}

\section{Strong solutions}

Consider
\begin{equation}\label{eq:sde}
\dd X_t=b(t,X_t)\dd t+\sigma(t,X_t)\dd W_t,\qquad X_0=x.
\end{equation}
A \emph{strong solution} is a continuous adapted process (adapted to the filtration generated by $W$) satisfying the integrated equation a.s. It is a machine: feed in the noise, get out the path.

\begin{theorem}[Existence and uniqueness]\label{thm:sdeexist}
Suppose $b,\sigma$ are Lipschitz in $x$ uniformly in $t$, and $|b(t,x)|+|\sigma(t,x)|\le K(1+|x|)$. Then \eqref{eq:sde} has a unique strong solution, with $\E\sup_{t\le T}|X_t|^2<\infty$.
\end{theorem}
\begin{proof}
\emph{Picard iteration.} Let $X^0\equiv x$ and $X^{n+1}_t=x+\int_0^tb(s,X^n_s)\dd s+\int_0^t\sigma(s,X^n_s)\dd W_s$. Using $(a+b)^2\le2a^2+2b^2$, Cauchy--Schwarz for the $\dd s$ integral, Doob and the isometry for the $\dd W$ integral, and Lipschitz continuity with constant $L$,
\[
\E\sup_{s\le t}|X^{n+1}_s-X^n_s|^2\le C\int_0^t\E\sup_{r\le s}|X^n_r-X^{n-1}_r|^2\dd s,\qquad C=2L^2(T+4).
\]
Iterating gives a bound $A(Ct)^n/n!$. Chebyshev and Borel--Cantelli give a.s.\ uniform convergence to a limit, which solves the equation.

\emph{Uniqueness.} If $X,Y$ are two solutions, let $v(t)=\E\sup_{s\le t\wedge\tau}|X_s-Y_s|^2$, stopped so as to be finite. The same estimate gives $v(t)\le C\int_0^tv(s)\dd s$. Gronwall's lemma then yields $v\equiv0$.
\end{proof}

\begin{lemma}[Gronwall]
If $v$ is bounded, nonnegative and $v(t)\le a+C\int_0^tv(s)\dd s$, then $v(t)\le ae^{Ct}$.
\end{lemma}

\paragraph{What goes wrong without the conditions.}
\begin{itemize}
\item \emph{Explosion.} $\dd X=X^2\dd t$, $X_0=1$ has solution $1/(1-t)$, which blows up at $t=1$. Superlinear growth lets the solution reach infinity. (Superlinear $\sigma$ does not always explode; it can be saved by negative drift, but the moments may be infinite.)
\item \emph{Non-uniqueness.} $\dd X=3X^{2/3}\dd t$, $X_0=0$ has both $X\equiv0$ and $X_t=t^3$.
\item \emph{No strong solution.} Tanaka's equation $\dd X=\sgn(X)\dd W$ has solutions (take $X$ a Brownian motion and define $W=\int\sgn(X)\dd X$, again Brownian). This is a \emph{weak} solution: the pair $(X,W)$ exists, but $X$ is not a function of $W$. Since $W=\int\sgn(X)\dd X$ can be recovered from $|X|$ by Tanaka's formula, $W$ loses the sign information, so no strong solution exists.
\end{itemize}

\begin{theorem}[Yamada--Watanabe, one dimension, stated]\label{thm:yw}
Pathwise uniqueness holds if $b$ is Lipschitz and $|\sigma(x)-\sigma(y)|\le c|x-y|^{1/2}$. Weak existence together with pathwise uniqueness implies strong existence.
\end{theorem}

This matters for CIR and Heston, where $\sigma(x)=\xi\sqrt x$ is not Lipschitz at zero.

\section{Solvable models}

\paragraph{Linear SDEs.} For $\dd X=(a_t+b_tX)\dd t+(c_t+e_tX)\dd W$, let $U_t$ be the homogeneous solution, $U_t=\exp(\int_0^t(b_s-\tfrac12e_s^2)\dd s+\int_0^te_s\dd W_s)$. Then
\[
X_t=U_t\Bigl(x+\int_0^t\frac{a_s-e_sc_s}{U_s}\dd s+\int_0^t\frac{c_s}{U_s}\dd W_s\Bigr),
\]
which one checks by applying the product rule to the claimed solution.

\paragraph{Ornstein--Uhlenbeck.} $\dd X=\kappa(\theta-X)\dd t+\sigma\dd W$. Multiply by $e^{\kappa t}$:
\begin{equation}\label{eq:ou}
X_t=\theta+(x-\theta)e^{-\kappa t}+\sigma\int_0^te^{-\kappa(t-s)}\dd W_s.
\end{equation}
So $X_t$ is Gaussian with mean $\theta+(x-\theta)e^{-\kappa t}$ and variance $\frac{\sigma^2}{2\kappa}(1-e^{-2\kappa t})$. Its stationary law is $\Nc(\theta,\sigma^2/2\kappa)$, and the half-life of a shock is $\ln2/\kappa$. Sampled discretely with step $h$, OU is an AR(1) with coefficient $e^{-\kappa h}$. This is the model behind Vasicek rates and pairs trading.

\paragraph{Cox--Ingersoll--Ross.} $\dd X=\kappa(\theta-X)\dd t+\sigma\sqrt{X}\dd W$. Taking expectations in the integrated equation gives $\E X_t=\theta+(x-\theta)e^{-\kappa t}$, and It\^o on $X^2$ gives
\[
\Var X_t=x\frac{\sigma^2}{\kappa}\bigl(e^{-\kappa t}-e^{-2\kappa t}\bigr)+\frac{\theta\sigma^2}{2\kappa}\bigl(1-e^{-\kappa t}\bigr)^2.
\]
The law is a scaled noncentral chi-squared: $X_t\sim c\,\chi'^2(\delta,\lambda)$, with $c=\sigma^2(1-e^{-\kappa t})/4\kappa$, degrees of freedom $\delta=4\kappa\theta/\sigma^2$, and noncentrality $\lambda=xe^{-\kappa t}/c$. The stationary law is $\mathrm{Gamma}(2\kappa\theta/\sigma^2,\ 2\kappa/\sigma^2)$ (shape, rate).

\begin{proposition}[Feller condition]\label{prop:feller}
If $2\kappa\theta\ge\sigma^2$, $X$ never reaches $0$ (from $x>0$). If $2\kappa\theta<\sigma^2$, it reaches $0$ a.s.\ and is instantaneously reflected.
\end{proposition}
This follows from the scale function criterion below. The dimension $\delta=4\kappa\theta/\sigma^2$ is the analogue of the dimension of a squared Bessel process, for which zero is hit exactly when $\delta<2$. In equity-calibrated Heston models, the Feller condition is often violated.

\section{The generator, Dynkin and exit problems}

Let $X$ solve a time-homogeneous version of \eqref{eq:sde}. Define the \emph{generator}
\[
\mathcal{A}f(x)=b(x)f'(x)+\tfrac12\sigma^2(x)f''(x).
\]
By It\^o's formula, $f(X_t)-\int_0^t\mathcal{A}f(X_s)\dd s$ is a local martingale. In particular, for a stopping time $\tau$ with enough integrability, we have \emph{Dynkin's formula}:
\[
\E_x f(X_\tau)=f(x)+\E_x\int_0^\tau\mathcal{A}f(X_s)\dd s.
\]

\begin{intuition}
$\mathcal{A}f(x)$ is the expected rate of change of $f(X)$ when $X$ sits at $x$: $\lim_{h\downarrow0}\frac{\E_xf(X_h)-f(x)}{h}$. It is the drift of $f(X_t)$. Every PDE in this chapter says one thing: \emph{the drift of something is zero (or equals a known running cost)}.
\end{intuition}

\paragraph{Exit problems.} Let $\tau$ be the exit time from $(a,b)$. Solving $\mathcal{A}s=0$ gives the \emph{scale function} $s'(x)=\exp(-\int^x2b/\sigma^2)$, and $s(X_t)$ is a local martingale. Therefore
\[
\Prob_x(X_\tau=b)=\frac{s(x)-s(a)}{s(b)-s(a)}.
\]
Solving $\mathcal{A}m=-1$ with $m(a)=m(b)=0$ gives the expected exit time $m(x)=\E_x\tau$. Boundary behaviour is read off $s$: a point at $0$ is not reached from $x>0$ if $s(0+)=-\infty$. For CIR, $s'(x)\propto x^{-2\kappa\theta/\sigma^2}e^{2\kappa x/\sigma^2}$, whose integral diverges at $0$ exactly when $2\kappa\theta/\sigma^2\ge1$. This proves Proposition~\ref{prop:feller} (up to behaviour at $\infty$).

\section{Kolmogorov equations}

\begin{theorem}[Backward equation]\label{thm:backward}
If $u\in C^{1,2}$ with suitable growth solves $\partial_tu+\mathcal{A}u=0$, $u(T,x)=h(x)$, then $u(t,x)=\E[h(X_T)\mid X_t=x]$.
\end{theorem}
\begin{proof}
By It\^o's formula, $u(s,X_s)$ has no drift; if it is a true martingale, then $u(t,X_t)=\E[u(T,X_T)\mid\F_t]=\E[h(X_T)\mid\F_t]$.
\end{proof}

\paragraph{Forward (Fokker--Planck) equation.} Let $p(T,y)$ be the density of $X_T$ given $X_t=x$. For any test function $f$, Dynkin's formula gives $\partial_T\E f(X_T)=\E\mathcal{A}f(X_T)$, that is, $\partial_T\int fp=\int(\mathcal{A}f)p$. Integrating by parts in $y$,
\begin{equation}\label{eq:fp}
\partial_Tp=-\partial_y(bp)+\tfrac12\partial_{yy}(\sigma^2p).
\end{equation}
The forward equation evolves \emph{probabilities} forward in time, and the backward equation evolves \emph{payoffs} backward. They are adjoint. Dupire's formula is an immediate consequence of the forward equation (Chapter~\ref{ch:localvol}).

Setting $\partial_T p=0$ gives the stationary density, for one-dimensional diffusion, $p_\infty(y)\propto\frac{1}{\sigma^2(y)}\exp(\int^y2b/\sigma^2)$ (speed density). Check for OU: this gives $\Nc(\theta,\sigma^2/2\kappa)$. For CIR it gives the stated Gamma.

\section{Feynman--Kac}

\begin{theorem}[Feynman--Kac]\label{thm:fk}
Let $u\in C^{1,2}([0,T)\times\R)\cap C([0,T]\times\R)$ solve
\[
\partial_tu+\mathcal{A}u-r(t,x)u+g(t,x)=0,\qquad u(T,x)=h(x),
\]
with $r$ bounded below. Suppose in addition that the process $Y_s=e^{-\int_t^sr}u(s,X_s)+\int_t^se^{-\int_t^vr}g(v,X_v)\dd v$ is a true martingale (for instance $\partial_xu\,\sigma$ with polynomial growth and polynomial moments). Then
\[
u(t,x)=\E_{t,x}\Bigl[e^{-\int_t^Tr(s,X_s)\dd s}h(X_T)+\int_t^Te^{-\int_t^sr(v,X_v)\dd v}g(s,X_s)\dd s\Bigr].
\]
\end{theorem}
\begin{proof}
By It\^o's formula and the PDE,
\[
\dd Y_s=e^{-\int_t^sr}\bigl(\partial_tu+\mathcal{A}u-ru+g\bigr)\dd s+e^{-\int_t^sr}\partial_xu\,\sigma\,\dd W_s=e^{-\int_t^sr}\partial_xu\,\sigma\,\dd W_s.
\]
So $Y$ is a local martingale, and under the assumption a martingale. Taking $\E_{t,x}$ of $Y_t=\E[Y_T\mid\F_t]$ gives the formula.
\end{proof}

\begin{pitfall}
The PDE alone does not determine the solution; you need growth conditions. The heat equation has non-zero solutions with zero initial data (Tychonoff's examples) that grow like $e^{x^{2+\epsilon}}$. Feynman--Kac selects the solution with controlled growth. In finance the same issue appears when the stock is a strict local martingale: the Black--Scholes-type PDE then has several solutions, and the price (the expectation) may not be the one with the ``expected'' boundary behaviour.
\end{pitfall}

\paragraph{The converse.} The class notes argue the other way: define $u$ by the expectation; then $u(s,X_s)$ is a martingale by the Markov property and the tower law, and if $u$ is smooth enough, its $\dd t$ term must vanish, which gives the PDE. Smoothness is the delicate part. It holds when $\sigma$ is uniformly elliptic and coefficients are H\"older.

\begin{example}[Black--Scholes both ways]
Under $\Q$, $\mathcal{A}=rx\partial_x+\tfrac12\sigma^2x^2\partial_{xx}$, so with $g=0$ the PDE is the BSM equation and the expectation is the risk-neutral price. The equation is thus the \emph{infinitesimal} form of $e^{-rt}V_t$ having zero drift.
\end{example}

\begin{example}[Bond prices]
In a short-rate model, $P(t,T)=\E_t[e^{-\int_t^Tr_s\dd s}]$ is given by Feynman--Kac with potential $r(x)=x$, $h=1$. This is the term-structure PDE of Chapter~\ref{ch:rates}.
\end{example}

\begin{example}[Occupation time / running payoffs]
$u(t,x)=\E_{t,x}[\exp(-\lambda\int_t^T\ind_{\{W_s>0\}}\dd s)]$ solves a PDE with a discontinuous potential. Solving it (Kac's original use) gives the arcsine law. In the same way, Parisian options and corridor products are priced with a discontinuous killing rate.
\end{example}

\section{Optimal stopping and American options}\label{sec:american}

An American option with payoff $g$ has value, under $\Q$,
\[
V(t,x)=\sup_{\tau\in[t,T]}\E^\Q_{t,x}\bigl[e^{-r(\tau-t)}g(S_\tau)\bigr].
\]
This is the Snell envelope (Theorem~\ref{thm:snell}) in continuous time. The discounted value is the smallest $\Q$-supermartingale above the discounted payoff. Therefore the \emph{drift of $e^{-rt}V$ is $\le0$}, and it is $0$ wherever it is optimal to continue. In formulas, with $\mathcal{L}=\partial_t+rx\partial_x+\tfrac12\sigma^2x^2\partial_{xx}-r$,
\begin{equation}\label{eq:vi}
\min\bigl(-\mathcal{L}V,\ V-g\bigr)=0,\qquad V(T,x)=g(x).
\end{equation}
This \emph{variational inequality} splits space into a continuation region $\{V>g\}$, where the BSM equation holds, and an exercise region $\{V=g\}$. The boundary between them is free, part of the solution.

\begin{theorem}[Early exercise of a call]
If $r\ge0$ and the stock pays no dividends, an American call is worth the European call. It is never optimal to exercise early.
\end{theorem}
\begin{proof}
Let $\tau\le T$ be a stopping time. By conditional Jensen, the martingale property of $e^{-rt}S_t$, and $e^{-rT}\le e^{-r\tau}$,
\[
\E^\Q[e^{-rT}(S_T-K)^+\mid\F_\tau]\ge(e^{-r\tau}S_\tau-Ke^{-rT})^+\ge e^{-r\tau}(S_\tau-K)^+.
\]
Taking expectations shows that holding to maturity dominates stopping at any $\tau$.
\end{proof}

\begin{intuition}
Exercising a call early pays the strike sooner (losing interest) and gives up the insurance against the stock falling below $K$. Both effects favour waiting. A put is different: exercising early \emph{receives} $K$ sooner, so with $r>0$ a deep in-the-money put should be exercised once the interest on $K$ outweighs the remaining insurance value.
\end{intuition}

\begin{theorem}[Perpetual American put]\label{thm:perp}
With $T=\infty$, $r>0$, the optimal rule is to exercise on first reaching
\[
L^*=\frac{2r}{2r+\sigma^2}K,
\]
and for $x\ge L^*$ the value is $V(x)=(K-L^*)\,(x/L^*)^{-2r/\sigma^2}$. For $x\le L^*$, $V(x)=K-x$.
\end{theorem}
\begin{proof}
\emph{Candidate.} In continuation, time does not matter, so $\tfrac12\sigma^2x^2V''+rxV'-rV=0$. Trying $x^p$ gives $p=1$ or $p=-\gamma$, where $\gamma=2r/\sigma^2$. As $V$ must stay bounded as $x\to\infty$, $V=Ax^{-\gamma}$. Value matching at $L$ gives $AL^{-\gamma}=K-L$. So for each threshold $L$ the value of the rule ``exercise at $L$'' is $V_L(x)=(K-L)(x/L)^{-\gamma}$. (It is also $(K-L)\E^\Q_x[e^{-r\tau_L}]$, and the Laplace transform of the hitting time gives this.) Maximising over $L$ gives $L^*=\gamma K/(1+\gamma)$.

\emph{Smooth pasting.} At $L^*$, $V'(L^*)=-\gamma(K-L^*)/L^*=-1$, matching the slope of $K-x$.

\emph{Verification.} $V\ge(K-x)^+$; $\mathcal{L}V=0$ on $x>L^*$; and for $x<L^*$, $\mathcal{L}V=-rK<0$. Hence $e^{-rt}V(S_t)$ is a supermartingale ($V$ being $C^1$ with a piecewise $C^2$ form, It\^o's formula still holds), and it dominates the discounted payoff, giving $V\ge$ every rule's value. Stopped at $\tau_{L^*}$ it is a martingale, so $V$ is attained.
\end{proof}

\begin{numcheck}
For $K=100$, $r=5\%$, $\sigma=20\%$: $L^*=71.43$ and $V(100)=12.3200$. A CRR tree at $T=100$ years with 20{,}000 steps gives $12.3179$. For a $T=1$ put with $S_0=100$, a tree with 5{,}000 steps gives $6.0902$, against the European $5.5731$: an early exercise premium of $0.5171$.
\end{numcheck}

\begin{exercises}
\item Solve $\dd X=-\frac{X}{T-t}\dd t+\dd W$ on $[0,T)$ with $X_0=0$, and show it is a Brownian bridge from $0$ to $0$.
\item For OU, compute $\Cov(X_s,X_t)$, and then $\Var\int_0^tX_s\dd s$. This variance appears in Vasicek bond prices.
\item For $\dd X=\mu X\dd t+\sigma X\dd W$, $\mu\ne\tfrac12\sigma^2$, find the scale function. What is the probability that it reaches $b$ before $a$?
\item Use Fokker--Planck to verify that the lognormal density solves the forward equation for GBM.
\item Show that a perpetual American call on a dividend-paying stock with yield $q>0$ has threshold $L^*=\frac{p}{p-1}K$, where $p>1$ solves $\tfrac12\sigma^2p(p-1)+(r-q)p-r=0$.
\item Let $h=(x-K)^+$. Compute $\partial_{xx}$ of the BSM solution near $T$, and show the PDE's smoothing: gamma blows up like $(T-t)^{-1/2}$ at $x=K$.
\item (Interview) Let $u(x)=\E_x[\int_0^\tau e^{-\lambda t}\dd t]$ where $\tau$ is the exit of Brownian motion from $(-1,1)$. Set up and solve the ODE.
\end{exercises}
